% 来源：OEBPS/Text/01_19.xhtml（OPF spine 顺序；仅第一册）
\cleardoublepage
\BookEntry[贾雷德·戴蒙德《第三种黑猩猩》]{贾雷德·戴蒙德《第三种黑猩猩》}

\BookCover{images/书单05第三种黑猩猩.jpg}{[美] 贾雷德·戴蒙德著| 王道还译| 上海译文出版社2016年版}

\begin{BookSummaryBox}
\textbf{人类作为一个物种数百万年来的演化史}
\par
\end{BookSummaryBox}

\subsubsection*{这本书谈什么？}

在自然界，人类与猩猩的关系最近，基因的差异不到2\%。人类与猩猩应属同一类目，黑猩猩只有两种，因此称人类是第三种猩猩也不为过。
\par

但人类如何变成世界的征服者，进而霸占这个世界绝大部分的资源？本书是人类这个物种数百年来的自然演化史。作者检视人类进化的过程，从语言、性行为和生命周期的角度，探讨人类与其他动物何以不同。全书充满令人好奇的问题，让读者重新认识自以为了解的人类。
\par

\subsubsection*{芒格为什么谈起这本书？}

在伯克希尔哈撒韦公司1999财年的年会上，芒格说：
\par

\begin{quotation}
我推荐大家去了解戴蒙德关于进化的观点。人类（第三种黑猩猩）是独一无二的，但是在艺术、语言、种族屠杀等方面，又和其他物种有相似之处。戴蒙德指出人类毁灭环境的倾向可能会让人类进化趋势逆转。他的警告发人深省。
\par
\end{quotation}

\subsubsection*{《第三种黑猩猩》作者贾雷德·戴蒙德如是说}

“大脑与直立姿势的演化，是语言与艺术的先决条件，但不是充分条件。骨骼形态像人，并不保证就有人性。我们要攀上人性的高峰，还得在生命循环上做彻底的改变。”
\par

“‘生命循环’指物种特定的生物特质，例如每胎生产的子女数目、亲职行为、成年个体间的社会关系、两性关系、两性互动的模式、性关系的频率、停经（如果有的话），以及平均寿命。我们认为人类的这些特质都是理所当然的，从未怀疑它们需要解释。但是我们的生命循环，要是以动物的常规来衡量，却是离谱的、奇异的。”
\par

“我们的祖先自从与其他的猿分化之后，就在几个重要的方面发生了变化，例如行进的体态与大脑的容量。而大猩猩与黑猩猩却没有发生过什么重大变化，与当年的人—猿共祖，模样没什么大差别。如果那是实情，人类可能与大猩猩最亲近，或与黑猩猩最接近，或者三者亲近的程度一样，彼此的遗传距离一样。”
\par
